% MAS1004 Assignment 1: long report.
%
% You fill this in with your agent, in your repository. AGENTS.md tells the
% agent how. Red text in the PDF is what still has to be replaced.
%
% - Pictures go in figures/. Copy them from results/, which is not in git.
% - Outputs go in outputs/. run.py saves what it prints in
%   results/<tag>_output.txt, and PLAN.md saves the other outputs in results/
%   with `| tee results/...`. Copy them here. \outputfile{outputs/check.txt}
%   prints a file exactly as it is. Every number in this report has to come
%   from such an output.
%
% To make the PDF: compress this folder into a zip, and on Overleaf choose
% New Project, Upload Project. Set Menu, Compiler to XeLaTeX, and Recompile.

\documentclass[11pt]{article}
\usepackage{mas1004}

\reporttitle{Assignment 1: Long report}
\studentname{김규화}
\studentid{20261167}

\begin{document}
\makeheader

% The pictures in figures/ and the numbers in this text were printed by the
% commands whose output is in outputs/, and by src/report_assets.py, which
% only copies and composes pictures and measures the saved models again.

% ---------------------------------------------------------------------------
\section*{Problem 1. Categories and images}

\subsection*{Why these categories}

Five kinds of cymbal from a drum kit: \texttt{crash\_cymbal},
\texttt{ride\_cymbal}, \texttt{hihat\_cymbal}, \texttt{splash\_cymbal} and
\texttt{china\_cymbal}. I am the drummer in my university band, so drums are
what I know best and what I am happy to look at for a week. The assignment asks
for categories that are hard to tell apart, and cymbals are: five round metal
plates that are separated by their size, by how far the edge is turned up or
down, by the bell in the middle, and by what is printed on the metal. The five I
chose are the five I actually play.

\subsection*{What I expected the model to use}

Before training I expected it to use shape: the size of the plate and how far
the edge is turned, because that is what I look at when I pick one up, and after
that the bell. I expected the colour and the background to be noise. I also
expected one class to fail. My five classes are five points on two scales, how
large the plate is and how far its edge is turned, and a crash is in the middle
of both: it is a plain, mostly flat plate, bigger than a splash and lighter than
a ride, so a model that reads those two scales has to put a crash in the middle
of them and answer \texttt{splash\_cymbal} or \texttt{ride\_cymbal} for one.

That is what happened, so the guess was right. In the run I published, of the
10 crash images in its 64-image test set 3 came out right and 5 were called
\texttt{splash\_cymbal} (\texttt{outputs/report\_numbers.txt}). On the 25 new
photographs it never once answered \texttt{crash\_cymbal} for anything, and the
5 crash photographs came out 0 right, 1 \texttt{china\_cymbal},
1 \texttt{hihat\_cymbal} and 3 \texttt{ride\_cymbal}
(\texttt{outputs/check.txt}); three of the five mistakes it was most sure about
are those crashes, at 75\%, 71\% and 50\% confidence. Problem 7 says why the
first half of the guess was wrong as well: the size of the plate and its turned
edge are not what it reads, it reads the middle of the frame first.

\subsection*{Images per category}

I asked the downloader for 120 images in each class. The middle column is what
it saved, and what \texttt{src/run.py} printed when it read the folder: for two
of the classes the search pages ran out before the 120th image. The right column
is what \texttt{src/clean.py count} left after cleaning, which Problem 4
describes.

\begin{center}\footnotesize
\begin{tabular}{lrr}
\toprule
class & images read by \texttt{run.py} & kept after cleaning \\
\midrule
\texttt{china\_cymbal}  & 82 & 70 \\
\texttt{crash\_cymbal}  & 60 & 51 \\
\texttt{hihat\_cymbal}  & 86 & 75 \\
\texttt{ride\_cymbal}   & 55 & 52 \\
\texttt{splash\_cymbal} & 80 & 76 \\
\midrule
all                     & 363 & 324 \\
\bottomrule
\end{tabular}
\end{center}

% ---------------------------------------------------------------------------
\section*{Problem 2. Preparing an image}

\sidebyside{figures/prepare_original.jpg}{\texttt{data/clean/splash\_cymbal/%
0001.jpg} as downloaded, $2018 \times 849$}
{figures/prepare_result.png}{the $224 \times 224$ that \texttt{prepare\_image}
returns, put back into ordinary colours}

\texttt{prepare\_image} keeps the shape and takes the shorter side to 256, so
this photo becomes $608 \times 256$, and then cuts the middle $224 \times 224$
out of it at $(192, 16)$. That throws away 192 pixels on the left, 192 on the
right, 16 at the top and 16 at the bottom: 32.2\% of the pixels the resize
left. In this photograph the whole cymbal is spread across the width of the
frame, so the left and the right edge of the cymbal are both outside the crop,
and what the model sees is the middle of the cymbal and the printing on it, not
the rim that would show how big it is. The red box below is that crop, drawn
back onto the photo as downloaded.

\begin{center}\picturefile[0.9\linewidth]{figures/prepare_crop_box.jpg}\end{center}

\texttt{ride\_cymbal/003.png} is the second image where this happens, and it is
one of the images from a new source in Problem 5. The cymbal fills the frame,
so the crop keeps the middle of it, the bell, the mounting nut and the printing,
and the rim, which is what tells a ride from a crash for me, is outside it. The
model answered \texttt{hihat\_cymbal} 86\%, the mistake it was most sure about
of all 25 new images.

\begin{center}\picturefile[0.55\linewidth]{figures/mistake_1.png}\end{center}

The two pictures above and the numbers in this paragraph come from
\texttt{src/report\_assets.py}. It prepares the same photo twice: once with
\texttt{prepare\_image}, and once with the steps \texttt{src/check.py} writes
again from \texttt{docs/model.json} for the web page. The two differ by
0.00000000, and \texttt{check.py} reports that the exported model agrees with
Python's answers (largest gap 0.0000), so the badge on the demo page is green.

% ---------------------------------------------------------------------------
\section*{Problem 3. First training runs}

Every row below is the line \texttt{run.py} printed at the end of that run,
copied with nothing changed. The whole output of each run is in \texttt{outputs/}.

\begin{center}\footnotesize
\begin{tabular}{llllrrrr}
\toprule
tag & start & trained & epochs & lr & trainable parameters & train acc. & test acc. \\
\midrule
run     & ImageNet & all layers & 10 & 0.0001 & 11,179,077 & 100.0\% & 62.5\% \\
scratch & random   & all layers & 10 & 0.0001 & 11,179,077 &  71.1\% & 37.5\% \\
frozen  & ImageNet & last layer & 10 & 0.001  &      2,565 &  75.3\% & 52.8\% \\
custom  & ImageNet & all layers & 15 & 0.0005 & 11,179,077 &  90.7\% & 50.0\% \\
clean   & ImageNet & all layers & 10 & 0.0001 & 11,179,077 &  99.2\% & 54.7\% \\
\bottomrule
\end{tabular}
\end{center}

\texttt{run} is the defaults, \texttt{scratch} is the same thing from random
numbers, \texttt{frozen} trains only the new last layer, and \texttt{custom} is
the setting I chose: 15 epochs at $5 \times 10^{-4}$ instead of the default 10
epochs at $1 \times 10^{-4}$. The last row is the run after cleaning, which
Problem 4 describes.

\picturefile{figures/run_curves.png}

The training loss falls from 1.3620 in epoch 1 to 0.0240 in epoch 10, and from
epoch 6 on it stays under 0.06: it flattens out. The test loss falls from 1.2723
to 1.1199 by epoch 5 and then climbs back to 1.2148, and the test accuracy
wanders between 58.3\% and 63.9\% after epoch 3, so it stopped improving long
before the training loss did.

The gap is 37.5 points: train accuracy is 100.0\% and test accuracy is 62.5\%.
With 291 training images the model can memorise them all, and the train loss
going down while the test loss goes up is what that memorising looks like.

\begin{center}\picturefile[0.5\linewidth]{figures/run_confusion.png}\end{center}

Measured again by \texttt{src/report\_assets.py}, the model gets 45 of the 72
test images right. The two it mixes up most are \texttt{crash\_cymbal} and
\texttt{ride\_cymbal}: 5 crashes are called rides and 3 rides are called
crashes. It also calls 3 splashes rides and 3 chinas hihats. The worst class is
\texttt{ride\_cymbal}, where only 5 of 11 are right, and the best is
\texttt{hihat\_cymbal} with 15 of 17.

Starting from ImageNet weights is worth 25.0 points here: 62.5\% against 37.5\%
from random numbers. The random model put 48 of the 72 test images into
\texttt{china\_cymbal} and never once answered \texttt{ride\_cymbal} or
\texttt{splash\_cymbal}, so with only 291 training images it learned the biggest
class and little else.

Training only the last layer comes close: 52.8\% against 62.5\%, 9.7 points
behind, with 2,565 of the 11,179,077 weights trainable and 0.8 minutes against
3.8. The ImageNet features already separate most cymbals, and the last layer on
its own is enough for 52.8\%. It leans on the same mistake, calling 5 crashes
rides, and adds its own: 4 splashes called crashes.

My longer run went the other way: 15 epochs at $5 \times 10^{-4}$ finished at
50.0\%, 12.5 points below the 10-epoch default, and it leaned on
\texttt{ride\_cymbal}, where it put 29 of the 72 test images and got only 7 of
the rides right. Training longer at a bigger step did not help.

\subsection*{The other runs}
\sidebyside{figures/scratch_curves.png}{scratch, from random numbers}
           {figures/frozen_curves.png}{frozen, last layer only}
\sidebyside{figures/custom_curves.png}{custom, 15 epochs at $5 \times 10^{-4}$}
           {figures/clean_curves.png}{clean, the run after cleaning (Problem 4)}
\sidebyside{figures/scratch_confusion.png}{\texttt{scratch}}
           {figures/frozen_confusion.png}{\texttt{frozen}}
\sidebyside{figures/custom_confusion.png}{\texttt{custom}}
           {figures/clean_confusion.png}{\texttt{clean}}

\subsection*{Output of the runs}
\outputfile{outputs/run_output.txt}
\outputfile{outputs/scratch_output.txt}
\outputfile{outputs/frozen_output.txt}
\outputfile{outputs/custom_output.txt}
\outputfile{outputs/clean_output.txt}

% ---------------------------------------------------------------------------
\section*{Problem 4. Cleaning}

\subsection*{My rule}

I kept a downloaded image only if it is an ordinary photograph of one whole
cymbal of that class: not wrapped in anything, not a close-up so tight that the
shape cannot be seen, and not a picture whose subject I cannot identify as a
cymbal at all. A photograph may appear only once. If the same photo, or a near
copy of it, had come back more than once, I kept one of them, the copy whose
class the photo really shows when the copies were in different classes, and
removed the others. The reasons recorded in \texttt{data/removed/log.csv} are
this rule written out one image at a time.

\subsection*{Three kinds of junk}

All 39 images I removed fall into three kinds. The pictures are the removed
files themselves, from \texttt{data/removed/}.

\sidebyside{figures/junk_1_different_class.jpg}{\texttt{china\_cymbal/0009.jpg}:
one \texttt{splash\_cymbal} product photo that the search also saved as
\texttt{china\_cymbal}, \texttt{crash\_cymbal} and \texttt{hihat\_cymbal} (7
images, 7 in all)}
{figures/junk_2_copy_same_class.jpg}{\texttt{china\_cymbal/0062.jpg}: a near
copy of \texttt{china\_cymbal/0015.jpg} in the same class (29 images)}

\begin{center}\picturefile[0.45\linewidth]{figures/junk_3_rule.jpg}\end{center}

\texttt{china\_cymbal/0048.jpg}: a cymbal wrapped in plastic, with the shape
barely visible. This is the third kind, images that break my rule: one extreme
close-up, one photo whose object is not identifiable as a cymbal, and this one.
Three images in all.

\subsection*{What I removed}
\outputfile{outputs/clean_count.txt}

So 39 of the 363 images went to \texttt{data/removed/}: 36 duplicates and 3 that
break the rule, and 324 were left for training.

\subsection*{The near copies}

The suspect search before cleaning listed 47 pairs of near copies
(\texttt{results/cleaning/suspects\_before/suspects.csv}). Every one of the 47
pairs lost at least one of its two files: in 33 pairs one of the two went, in 14
pairs both went, in none did both stay. It cannot print what it found with
\texttt{suspects.csv} alone, so the pair rows at the end of the file are listed
in \texttt{outputs/report\_numbers.txt} with, for each pair, which of the two I
removed. 20 of the removals record how alike the two were, from 0.9516 to
0.9998. After cleaning I ran the same search again and it found no near copies
left, and the agreement between the label and the rest of the data, which was
59.2\% over 363 images, is 57.4\% over 324. Read that as: the copy that stayed is
still a little unusual, but the copy itself is gone.

\outputfile{outputs/suspects.txt}
\outputfile{outputs/suspects_after.txt}

\subsection*{The first 20 suspects in each class}

\begin{center}\footnotesize
\begin{tabular}{lrrr}
\toprule
class & in the first 20 & removed & kept \\
\midrule
\texttt{china\_cymbal}  & 20 & 4 & 16 \\
\texttt{crash\_cymbal}  & 20 & 3 & 17 \\
\texttt{hihat\_cymbal}  & 20 & 2 & 18 \\
\texttt{ride\_cymbal}   & 20 & 0 & 20 \\
\texttt{splash\_cymbal} & 20 & 0 & 20 \\
\midrule
all five                & 100 & 9 & 91 \\
\bottomrule
\end{tabular}
\end{center}

Of the first 20 suspects in each class I removed 9 of the 100. They are
\texttt{china\_cymbal} ranks 3, 6, 9 and 20, \texttt{crash\_cymbal} ranks 9, 18
and 20, and \texttt{hihat\_cymbal} ranks 4 and 14; the \texttt{ride\_cymbal} and
\texttt{splash\_cymbal} lists held nothing I wanted to remove, which is what
that list is for: it shows the images that look least like the rest of their
class, and many of those are good pictures that are only unusual. All three
images that break my rule are in the \texttt{china\_cymbal} list, at ranks 3, 6
and 9, and the rest are copies, which is the kind of junk I did not find by
looking at 363 images one by one and the reason the search is worth running.

\subsection*{Before and after cleaning}
\outputfile{outputs/compare.txt}

On the 25 new images, which cleaning does not touch at all, the same two saved
models go from 36.0\% to 52.0\%: cleaning is worth 16.0 points there, and that is
the one fair comparison in this section. The test accuracy goes from 62.5\% to
54.7\%, but the two numbers are measured on different test sets: the first run
tested on 72 images and 7 of those 72 are in \texttt{data/removed/} now, and the
cleaned test set of 64 images has no near copies left in it, which makes it the
harder set of the two. The cleaned model also learns its training set less
exactly, 99.2\% against 100.0\%, now that the copies that made it easy are gone.

\picturefile{figures/clean_worst.png}

The ten mistakes the cleaned model was most sure about, on its own test set, with
what the model said printed on each picture (\texttt{src/evaluate.py} draws it).
The worst is a splash it called \texttt{hihat\_cymbal} with 98\%, then two rides it
called \texttt{crash\_cymbal} with 97\% and 94\%, and a splash and a crash it
called \texttt{ride\_cymbal} with 79\% and 76\%.

% ---------------------------------------------------------------------------
\section*{Problem 5. Images from a new source}

The 25 images came from two places, and neither of them is the search the
download worked from. Four images in each class are the product photographs the
makers and the shops publish of the exact models I play, saved one at a time
from their own pages, and the fifth is my own photograph of the cymbal in front
of me, taken with my phone in the rehearsal room (1050 by 1400 pixels, or 1050
by 564 for one china). None of them is in the download: the overlap check
compares all 25 with all 363 downloaded images and finds no near copy, and the
closest any of them comes is 0.89, where a near copy counts above 0.95.

Five images for each class, 25 in all, in \texttt{data/new\_images/}, saved as
JPEG and PNG. The overlap check compares all 25 with all 363 downloads.

\outputfile{outputs/overlap.txt}

The same 25 images are also the only fair way to compare the model before and
after cleaning, because cleaning never touches them:

\begin{center}\footnotesize
\begin{tabular}{lrr}
\toprule
run & downloaded test set & the 25 new images \\
\midrule
run, before cleaning    & 62.5\% (72 images) & 36.0\% (9 of 25) \\
clean, after cleaning   & 54.7\% (64 images) & 52.0\% (13 of 25) \\
\bottomrule
\end{tabular}
\end{center}

Both rows are what \texttt{check.py --compare run clean} printed, in
\texttt{outputs/compare.txt}. On my own images the model gains 16.0 points from
cleaning. On the downloaded test set the number looks 7.8 points worse, but
that set is not the same set any more, as Problem 4 explains.

\subsection*{The five mistakes it was most sure about}

These are the five printed at the end of \texttt{outputs/check.txt}, measured by
check.py on the exported model. \texttt{src/report\_assets.py} found the same
five, and copied the pictures out of \texttt{data/new\_images/}.

\sidebyside{figures/mistake_1.png}{\texttt{ride\_cymbal/003.png}: said
\texttt{hihat\_cymbal} (86\%)}
{figures/mistake_2.png}{\texttt{crash\_cymbal/003.png}: said
\texttt{hihat\_cymbal} (75\%)}

\sidebyside{figures/mistake_3.png}{\texttt{crash\_cymbal/001.png}: said
\texttt{ride\_cymbal} (71\%)}
{figures/mistake_4.jpg}{\texttt{crash\_cymbal/005.jpg}: said
\texttt{ride\_cymbal} (50\%)}

\begin{center}\picturefile[0.5\linewidth]{figures/mistake_5.jpg}\end{center}

\texttt{splash\_cymbal/005.jpg}: said \texttt{china\_cymbal} (50\%).

\begin{itemize}
\item \texttt{ride\_cymbal/003.png}: a Paiste Monad Dry Heavy Ride of 22 inches,
photographed from above so that the cymbal fills the picture. The frame cuts
the rim off all around it, so the square the model gets holds the bell, the
mounting nut and the large printing, on a dark background with a piece of the
kit at the lower left. The cymbal is purple, and I have no purple cymbal in my
downloads.
\item \texttt{crash\_cymbal/003.png}: a Sabian SBR cymbal from above, with a
black felt, the wing nut and the stand through the middle. The cymbal again
fills the frame and its edge is outside the picture, and the \texttt{SABIAN}
printing is in the middle of the square, on a blurred room.
\item \texttt{crash\_cymbal/001.png}: a very small copy of the same Sabian XHH
cymbal, 169 by 102 pixels, the smallest file I collected. It shows only the
middle of the cymbal: the hole, the red plastic sleeve on the stand and the
\texttt{XHH} printing.
\item \texttt{crash\_cymbal/005.jpg}: a Meinl HCS 16" crash, photographed from
directly above on the kit. The cymbal fills the picture, its rim is cut by the
frame at the top and the bottom, and other parts of the kit (a stand, another
cymbal, a drum) are at the edges. The answer was a 50/50 one, so this is a
mistake it was barely sure about.
\item \texttt{splash\_cymbal/005.jpg}: an Istanbul Agop XIST 10" Brilliant
Splash on its stand in a rehearsal room, from above, with the floor, a stand
and a case around it. The whole cymbal is inside the frame here, and the answer
was 50/50 again.
\end{itemize}

Two of the five were called \texttt{hihat\_cymbal} and two \texttt{ride\_cymbal};
the fifth was called \texttt{china\_cymbal}. In all five the square holds the
middle of the cymbal and not its rim, and in three of them the cymbal is
mounted with the stand through the middle of the frame.

\subsection*{What is different about these images}

\begin{itemize}
\item The cymbals are mounted. Every one of these five is a cymbal on a kit or
on a stand, and the mount is in the middle of the frame: a felt, a wing nut, a
clutch, a red plastic sleeve. In my downloads the cymbal is usually off the
kit or on a plain stand, so nothing is through the middle of it.
\item The room is in the picture: a wooden wall, a blue wall, a floor, cables,
mic stands, a case, other cymbals and drums. My downloads are almost all
catalogue photos, with the cymbal alone on white or on black.
\item The camera is close and high. From directly above or from the side, at
arm's length, the cymbal fills the frame and the camera cuts its rim off. The
224 square then holds the middle of the cymbal, the bell, the mount and the
printing, while the outline and the size, which the catalogue photos show
whole, are outside it.
\item The light is indoor and uneven, with one strong reflection across the
metal, while the catalogue photos are lit evenly from the front.
\item The files themselves are a mixture: the PNGs run from 169 by 102 to 1360
by 535 pixels, four of the five JPEGs are 1050 by 1400, and
\texttt{china\_cymbal/005.jpg} is 1050 by 564.
\end{itemize}

% ---------------------------------------------------------------------------
\section*{Problem 6. The demo}

Address: \url{https://gyuhwa2007-cloud.github.io/mas1004-assignment1/}

The page serves the cleaned model (\texttt{python src/export\_web.py --tag
clean}). Its three files, \texttt{docs/model.onnx} (44.8 MB),
\texttt{docs/model.json} and \texttt{docs/selftest.json}, are committed in the
repository and GitHub Pages publishes the \texttt{docs/} folder of the
\texttt{main} branch, so the address works without this computer. The badge at
the top is green: \texttt{check.py} reports that the exported model agrees with
Python's answers, with a largest gap of 0.0000, and in
\texttt{src/report\_assets.py} the web page's pipeline and
\texttt{prepare\_image} differ by 0.00000000.

One thing to read carefully there: what those two measurements compare is the
crop, the arithmetic and the network on a picture that is already 256 pixels
wide, which is what the three self test pictures are (\texttt{selftest\_picture}
in \texttt{src/export\_web.py}). The page shrinks a photograph you give it
itself, in JavaScript, by halving it step by step and letting the browser smooth
each step (\texttt{docs/app.js}), which lands close to the single bilinear
resize \texttt{prepare\_image} does but not on the same numbers. So the
percentages this report quotes are the ones Python printed, and the percentages
the page shows for a full-size photograph are its own; the class it answers is
the same one, unless two classes are almost tied.

The page has no `none of these' row: it always shows its five classes with a
percentage each, so something that is none of them still comes back as one of
them. Shown a photograph of a parked bicycle it answered \texttt{hihat\_cymbal}
42.6\%, and given a second bicycle it answered \texttt{hihat\_cymbal} 79.1\% and
was sure of it. Both numbers are in \texttt{outputs/problem7.txt}, measured on
the exported model through the same pipeline the page runs
(\texttt{src/problem7.py}), and both bicycles are Wikimedia Commons photographs
(CC BY 2.0 and CC BY-SA 4.0).

\sidebyside{figures/none_bicycle_1.jpg}{a bicycle under a street lamp, none of
my classes: the page says \texttt{hihat\_cymbal} 42.6\%}
{figures/none_bicycle_2.jpg}{a bicycle in Rome, none of my classes: the page
says \texttt{hihat\_cymbal} 79.1\%}

% ---------------------------------------------------------------------------
\section*{Problem 7. What the model looks at}

\subsection*{Guess 1: the model reads the name printed on the cymbal}

My guess, from the china that came back wrong on the page: the model reads the
name of the type printed on the metal. In the china image that failed the
printing is large and lies differently from the other images, so I covered it
and tried the same image again.

The change: a grey box over the printing, at the same place in the $224 \times
224$ square the model gets from every photograph (the boxes are fractions of
that square, so they land on the same part of the cymbal in each image). If the
printed name is what the model reads, the answer should move when it is covered;
if the answer does not move, the name is not what it uses.

My request to the agent, word for word (my own words for this experiment, from
my short report):
\begin{quote}
“종류 이름이 나와 있는 문구를 가리고 테스트하였다.” --- ``I covered the
writing that shows the name of the type, and tested it again.''
\end{quote}

\sidebyside{figures/photo_1.jpg}{\texttt{china\_cymbal/003.png} as it is:
\texttt{china\_cymbal} 48.6\%}
{figures/guess1_changed_1.jpg}{printing covered:
\texttt{ride\_cymbal} 32.8\%, \texttt{china\_cymbal} 30.8\%}

\sidebyside{figures/photo_2.jpg}{\texttt{crash\_cymbal/003.png} as it is:
\texttt{hihat\_cymbal} 75.2\%}
{figures/guess1_changed_2.jpg}{printing covered: \texttt{hihat\_cymbal} 76.3\%}

\sidebyside{figures/photo_3.jpg}{\texttt{splash\_cymbal/005.jpg} as it is:
\texttt{china\_cymbal} 50.3\%}
{figures/guess1_changed_3.jpg}{printing covered: \texttt{china\_cymbal} 47.8\%}

\sidebyside{figures/photo_4.jpg}{\texttt{hihat\_cymbal/005.jpg} as it is:
\texttt{hihat\_cymbal} 86.3\%}
{figures/guess1_changed_4.jpg}{printing covered: \texttt{hihat\_cymbal} 78.7\%}

\begin{center}\footnotesize
\begin{tabular}{llrr}
\toprule
photo & really & before & after \\
\midrule
1 & \texttt{china\_cymbal/003.png}  & \texttt{china} 48.6\% & \texttt{ride} 32.8\% \\
2 & \texttt{crash\_cymbal/003.png}  & \texttt{hihat} 75.2\% & \texttt{hihat} 76.3\% \\
3 & \texttt{splash\_cymbal/005.jpg} & \texttt{china} 50.3\% & \texttt{china} 47.8\% \\
4 & \texttt{hihat\_cymbal/005.jpg}  & \texttt{hihat} 86.3\% & \texttt{hihat} 78.7\% \\
\bottomrule
\end{tabular}
\end{center}

What the results say: on the china the cover moved the answer, from
\texttt{china\_cymbal} 48.6\% to \texttt{ride\_cymbal} 32.8\% with the china just
behind at 30.8\%, so the printing carries a large part of that answer and the
new answer is a near tie. On the other three photographs nothing moved: the
hihat lost 7.6 points and stayed the answer, and both mistakes stayed exactly
where they were, a crash still called \texttt{hihat\_cymbal} 76.3\% and a splash
still called \texttt{china\_cymbal} 47.8\%. One photograph out of four is not a
rule, so the guess is too strong: the model reads the printing sometimes, and
something else decides the other three.

\subsection*{Guess 2: the model uses the bell and the mount in the middle of the
frame}

My guess, from the hihat and the ride in my short report: what the model uses is
what is holding the cymbal, in the middle of the picture. A hihat, unlike the
others, is fixed to a particular stand, and a hihat away from its stand had much
lower accuracy; a ride is recognised by its large, clear bell, and the rides
without one fall behind. Covering the bell and the mount should then move the
answer more than anything else I try.

The change: one grey box over the middle of the $224 \times 224$ square, so that
it covers the bell, the mounting nut and the clutch or felt of the stand, on
five photographs.

My request to the agent, word for word (my own words for this experiment, from
my short report):
\begin{quote}
“하이햇은 다른 심벌들과 달리 특정 스탠드와 함께 붙어있는 큰 특징을 가지기.
따라서 스탠드 없이 있는 하이햇을 가져다 두었을 때에는 현저히 낮은 정확도를
보여주었다.” --- ``A hihat, unlike the other cymbals, has the big feature of
being fixed to a particular stand, so a hihat that is not on a stand had much
lower accuracy.''

“라이드 심벌은 모델이 특징으로 크기가 크고 뚜렷한 벨로 잡았다.” --- ``The
model takes the large, clear bell of a ride as its feature.''
\end{quote}

\sidebyside{figures/photo_1.jpg}{\texttt{china\_cymbal/003.png} as it is:
\texttt{china\_cymbal} 48.6\%}
{figures/guess2_changed_1.jpg}{bell and mount covered:
\texttt{splash\_cymbal} 35.8\%}

\sidebyside{figures/photo_2.jpg}{\texttt{crash\_cymbal/003.png} as it is:
\texttt{hihat\_cymbal} 75.2\%}
{figures/guess2_changed_2.jpg}{bell and mount covered:
\texttt{ride\_cymbal} 40.1\%}

\sidebyside{figures/photo_3.jpg}{\texttt{splash\_cymbal/005.jpg} as it is:
\texttt{china\_cymbal} 50.3\%}
{figures/guess2_changed_3.jpg}{bell and mount covered:
\texttt{china\_cymbal} 45.6\%}

\sidebyside{figures/photo_4.jpg}{\texttt{hihat\_cymbal/005.jpg} as it is:
\texttt{hihat\_cymbal} 86.3\%}
{figures/guess2_changed_4.jpg}{bell and mount covered:
\texttt{hihat\_cymbal} 45.5\%}

\sidebyside{figures/photo_5.jpg}{\texttt{ride\_cymbal/003.png} as it is:
\texttt{hihat\_cymbal} 85.8\%, the ride itself 3.8\%}
{figures/guess2_changed_5.jpg}{bell and mount covered:
\texttt{hihat\_cymbal} 33.3\%, the ride 18.9\%}

\begin{center}\footnotesize
\begin{tabular}{llrr}
\toprule
photo & really & before & after \\
\midrule
1 & \texttt{china\_cymbal/003.png}  & \texttt{china} 48.6\% & \texttt{splash} 35.8\% \\
2 & \texttt{crash\_cymbal/003.png}  & \texttt{hihat} 75.2\% & \texttt{ride} 40.1\% \\
3 & \texttt{splash\_cymbal/005.jpg} & \texttt{china} 50.3\% & \texttt{china} 45.6\% \\
4 & \texttt{hihat\_cymbal/005.jpg}  & \texttt{hihat} 86.3\% & \texttt{hihat} 45.5\% \\
5 & \texttt{ride\_cymbal/003.png}   & \texttt{hihat} 85.8\% & \texttt{hihat} 33.3\% \\
\bottomrule
\end{tabular}
\end{center}

What the results say: this is the change that matters. Two of the five answers
moved to another class, the china to \texttt{splash\_cymbal} 35.8\% and the
crash it had called a hihat to \texttt{ride\_cymbal} 40.1\%, and the two
photographs it had right lost more than half of their confidence, the hihat
going from 86.3\% to 45.5\% and the wrong \texttt{hihat\_cymbal} on the ride
from 85.8\% to 33.3\%. On that ride the cymbal's own class rose from 3.8\% to
18.9\% the moment the middle of the frame was covered, so what made the model
answer hihat about a ride was the bell and the mounting hardware, not the plate.
Covering the middle cost the model every answer it had and gave it no better
one. My hihat and ride guesses were right about what the model looks at, and
wrong about the rest of the cymbal not mattering: the model leans on the middle
of the frame more than on anything else.

\subsection*{Guess 3: the frame cuts the turned edge off, and that is the
problem}

My guess, from the china again: the model uses how far the edge of a china is
turned back, and the reason it fails is that the cymbal is photographed so close
that the edge, and with it the shape, are outside the square the model sees. My
short report put it this way:

\begin{quote}
“차이나 심벌들의 이미지를 보면 가장자리의 거꾸로 꺾인 정도가 심할수록 더 높은
확률을 가진다는 것을 알 수 있다. ... 그러나 모델이 보는 이미지가 확대되어
가장자리가 잘 안 보인다. 따라서 너무 확대된 사진을 본다는 이유로 잘못된 결과를
낸다고 볼 수 있다.” --- ``Looking at the china images, the more sharply the edge
is turned back the higher the probability, so the model reads that turned edge.
But the edge cannot be seen well because the image the model sees is zoomed in,
so the wrong answer can be put down to it seeing a photograph that is too zoomed
in.''
\end{quote}

The change: shrink each photograph onto a grey margin, so that the whole cymbal,
edge and outline included, fits inside the $224 \times 224$ square with the
cymbal smaller than it was. If the missing edge is the problem, the mistakes
should go away.

\sidebyside{figures/photo_1.jpg}{\texttt{china\_cymbal/003.png} as it is:
\texttt{china\_cymbal} 48.6\%}
{figures/guess3_changed_1.jpg}{whole china on a grey margin:
\texttt{splash\_cymbal} 49.3\%}

\sidebyside{figures/photo_2.jpg}{\texttt{crash\_cymbal/003.png} as it is:
\texttt{hihat\_cymbal} 75.2\%}
{figures/guess3_changed_2.jpg}{whole crash on a grey margin:
\texttt{hihat\_cymbal} 86.8\%}

\sidebyside{figures/photo_3.jpg}{\texttt{splash\_cymbal/005.jpg} as it is:
\texttt{china\_cymbal} 50.3\%}
{figures/guess3_changed_3.jpg}{whole splash on a grey margin:
\texttt{china\_cymbal} 32.0\%}

\sidebyside{figures/photo_4.jpg}{\texttt{hihat\_cymbal/005.jpg} as it is:
\texttt{hihat\_cymbal} 86.3\%}
{figures/guess3_changed_4.jpg}{whole hihat on a grey margin:
\texttt{hihat\_cymbal} 81.9\%}

\begin{center}\footnotesize
\begin{tabular}{llrr}
\toprule
photo & really & before & after \\
\midrule
1 & \texttt{china\_cymbal/003.png}  & \texttt{china} 48.6\% & \texttt{splash} 49.3\% \\
2 & \texttt{crash\_cymbal/003.png}  & \texttt{hihat} 75.2\% & \texttt{hihat} 86.8\% \\
3 & \texttt{splash\_cymbal/005.jpg} & \texttt{china} 50.3\% & \texttt{china} 32.0\% \\
4 & \texttt{hihat\_cymbal/005.jpg}  & \texttt{hihat} 86.3\% & \texttt{hihat} 81.9\% \\
\bottomrule
\end{tabular}
\end{center}

What the results say: showing the model the whole cymbal fixed nothing. The
crash it had called a hihat was called a hihat even more strongly, 86.8\%, and
the china left its own class for \texttt{splash\_cymbal} 49.3\%. A splash stayed
\texttt{china\_cymbal} 32.0\% with \texttt{hihat\_cymbal} 28.2\% close behind, so
it is still wrong but no longer sure of it. Not one of the four answers became
right. So the turned edge and the missing outline are not what the model is
reading: if they were, giving it the whole cymbal would have helped, and it hurt
instead.

\subsection*{What I now think the model uses}

\begin{itemize}
\item The middle of the frame first of all, whatever is in it: the bell, the
mounting nut, the clutch, the felt, the stand. Covering that one place (Guess 2)
is the only change that moved two answers to another class and halved the
confidence of the rest, and on the ride photograph it was the difference between
\texttt{hihat\_cymbal} 85.8\% and \texttt{ride\_cymbal} 18.9\%.
\item A hihat is the class it can do: 5 of 5 new images right, the only class it
never misses, and the class it answers when it is shown a bicycle (79.1\%) or a
ride from above (85.8\%).
\item A ride is thin evidence: 2 of 5, the rest called \texttt{hihat\_cymbal}
twice and \texttt{china\_cymbal} once. On the ride I covered, the cymbal's own
class went from 3.8\% to 18.9\% when the middle of the frame was hidden, so a
large, clear bell is not enough on its own.
\item A crash is the worst: 0 of 5, called \texttt{ride\_cymbal} three times.
What my short report said about the crash survives: it is the class with no big
feature of its own, so it is answered with whatever feature the model does see,
here the bell and the mounting hardware of other classes.
\item The printing on the metal counts, sometimes: covering it moved 1 of the 4
photographs, and it moved that one to \texttt{ride\_cymbal}, which is not the
name printed on the china.
\item The turned edge of a china, which I thought was the feature, is not what
the model needs: giving it the whole cymbal, edge included, made the mistakes
worse.
\end{itemize}

How sure I am: the middle of the frame is the one thing the three changes agree
about, and it is measured on five photographs, so I am confident about it. The
printing rests on one photograph, so I am not confident about it at all.
Everything above is measured on the 5 to 25 images I have, and no change is
measured on more than five photographs.

\subsection*{Output of the changes}
\outputfile{outputs/problem7.txt}

% ---------------------------------------------------------------------------
\section*{Appendix. Output of check.py}
\outputfile{outputs/check.txt}

Every number in this report is one that a file in \texttt{outputs/} printed,
or one that \texttt{src/report\_assets.py} printed into
\texttt{outputs/report\_numbers.txt} while measuring the saved models again on
the test sets they were trained on, or a difference between two of those
numbers, or a percentage that \texttt{src/evaluate.py} printed on one of the
figures in \texttt{figures/}. The pictures in \texttt{figures/} were
copied out of \texttt{results/} by \texttt{src/report\_assets.py}, except the
single-image ones, which it made from \texttt{data/clean/},
\texttt{data/removed/} and \texttt{data/new\_images/}, and except the ones for
Problem 7, which \texttt{src/problem7.py} made from \texttt{data/new\_images/}
and from the changed copies it wrote into \texttt{data/changed/}. Nothing in
\texttt{data/} was changed.

\end{document}